\section{证据}
\begin{enumerate}
  \item 3 大属性、3 大原则、4 大分类、9 大种类、7 大规则。33497
  \item 可以证明案件事实的材料（中性词） \yjt 查证属实 \yjt 证据。
  \item 3 大属性。同时满足。
        \begin{enumerate}
          \item 客观性/真实性。\al{证人与当事人的利害关系 \yjt 客观性}。
          \item 合法性。主体、形式、程序、质证。
          \item 关联性。与犯罪的动机、起因、方式、结果相关的任何客观联系。\al{猥亵儿童，手机里的色情照片 \dui}。
                \begin{enumerate}
                  \item 关联性不要求必然具有因果关系。
                  \item \al{一并提供曾犯制造毒品罪的裁判文书 \dui{累犯关联}}。
                  \item 类似行为。嫌疑人/被告人曾经实施过与被指控罪行类似的犯罪行为。\al{可用于证明主观目的 \cuo}。
                  \item 品格材料。证明嫌疑人/被告人品性道德状况的材料。
                \end{enumerate}
        \end{enumerate}
  \item 3 大原则。
        \begin{enumerate}
          \item 证据裁判。没有证据/证据不足，不得认定事实。
                \begin{enumerate}
                  \item 案件事实的认定必须要依靠证据。
                  \item 定案的证据必须要有证据的能力。
                  \item 经过法庭调查过的查证属实的证据。
                  \item 要满足证明的标准。
                \end{enumerate}
          \item 自由心证。同民诉。
          \item 直接言词。亲自面对，口头陈述。
        \end{enumerate}
  \item 4 大理论分类。
        \begin{enumerate}
          \item 原始证据、传来证据（复制、转述）。\al{手机提取的视频}。
          \item 有罪和无罪。
          \item 言词和实物。言词形式提供。\al{证人证言、被害人陈述、被告人供述和辩解；鉴定意见}。
          \item 直接和间接。单独、直接（无需推理）证明案件主要事实（犯罪行为是否系嫌疑人/被告人所实施）。
                \al{甲告知一声自己系被车辆撞倒 \cuo}。
                \begin{enumerate}
                  \item 直接证据 \bdy 单独定案的证据。
                        \begin{enumerate}
                          \item 肯定性直接证据。需要肯定 2 个，谁、犯罪行为。
                          \item 否定性直接证据。只需否定 1 个，不要求是否真实 \yjt 查证那是别人的事。
                                \al{甲辩称公诉人指控他洗钱的 1093 万都是他的合法所得，该辩解是直接证据 \dui}。
                        \end{enumerate}
                  \item 与证据的真伪、来源无关。
                  \item 没有直接证据，间接证据满足条件也可以定案。
                \end{enumerate}
        \end{enumerate}
  \item 9 大种类 + 审查判断规则。之外都是非法形式。\al{外国判决}。
        \begin{enumerate}
          \item 物证。物理属性、痕迹、外部特征。间接证据。\al{警服}。
          \item 书证。内容和思想。内容有关，物质载体。诉外形成。\sjt 真伪不明、来源不明、无法解释。
                \al{无法证明来源，遂要求侦查人员进行补正 \hz 作出合理解释 \cuo}。
                \begin{enumerate}
                  \item \al{侦查机关制作的失窃药材清单是书证 \cuo{侦查笔录}}。
                  \item \al{从通信公司调取的通话记录清单是书证 \dui{用通话内容查实销赃情况}}。
                  \item \al{存折是物证 \cuo{书证}}。
                \end{enumerate}
          \item 证人证言。
          \item 被害人陈述。
          \item 嫌疑人/被告人供述 + 辩解 = 口供。\sjt 看身份判断。
          \item 鉴定意见。书面意见。刑诉只有指派 \hz 聘请（保证客观公证）。
                \al{对毒品含量进行鉴定的鉴定人甲开庭前死亡，该鉴定意见不得作为定案依据 \cuo{有成为的可能}}。
          \item 勘验、检查、辨认、侦查实验等笔录（扩大解释，\al{拍照}）。客观如实的记录。
          \item 视听资料（模拟信号） \ydy 录音录像 vs 电子数据（电子信号）。诉外。真伪不明、无法解释。\al{被篡改 \dui}。
                \begin{enumerate}
                  \item 案发中拍的 \yjt 视听资料。
                  \item 传到网上 \yjt 电子数据。
                  \item 讯问录像 \yjt 口供。
                  \item 用于证明没有刑讯逼供 \yjt 视听资料。
                  \item 对电子数据介质的操作都是应当。\al{应将 U 盘拆封过程进行录像 \dui}。
                  \item 注意和对电子数据本身的提取区分。\al{提取电脑数据时应全程录像 \cuo{有条件的}}。
                \end{enumerate}
          \item 个人出具的专门性问题报告、有关部门出具的事故调查报告 \ydy 鉴定意见。\al{物价部门的价格认定报告}。
        \end{enumerate}
  \item 收集。
        \begin{enumerate}
          \item 公检法、辩护律师。
          \item 行政机关。实物证据可以直接用，言词证据见错就排（不是不能用）。\al{甲在行政执法时的笔录不能作为证据使用 \cuo}。
          \item 监察机关。+证人证言、口供。
          \item 境外证据。
        \end{enumerate}
  \item 移送。
        \begin{enumerate}
          \item 全部移送。未移送的，应当通知，应当根据在案证据对案件事实作出认定。
          \item 应当对讯问过程录音录像的案件，必要时可以通知移送。
        \end{enumerate}
\end{enumerate}
